\documentclass[UTF8,zihao=-4]{ctexart}

\usepackage[a4paper,margin=1.8cm]{geometry}
\usepackage{amsmath,amssymb}
\usepackage{array,tabularx,longtable,multirow}
\usepackage{enumitem}
\usepackage{listings}
\usepackage{hyperref}
\usepackage{fancyhdr}
\usepackage{lastpage}

\setCJKmainfont{Source Han Serif CN}
\setCJKsansfont{Source Han Sans CN}
\setCJKmonofont{Noto Sans Mono CJK SC}
\setmonofont{DejaVu Sans Mono}

\pagestyle{fancy}
\fancyhf{}
\cfoot{第 \thepage 页共 \pageref{LastPage} 页}
\renewcommand{\headrulewidth}{0pt}

\setlength{\parindent}{0pt}
\setlength{\parskip}{0.35em}
\renewcommand{\arraystretch}{1.25}
\setlist[enumerate]{leftmargin=2.2em,itemsep=0.45em,topsep=0.2em}

\lstset{
  basicstyle=\ttfamily\small,
  columns=fullflexible,
  keepspaces=true,
  showstringspaces=false,
  xleftmargin=1.5em
}

\newcommand{\score}[1]{\textbf{(#1 分)}}
\newcommand{\eps}{\varepsilon}
\newcommand{\gto}{\rightarrow}
\newcommand{\kw}[1]{\texttt{#1}}
\newcommand{\qtitle}[3]{\vspace{0.8em}\textbf{#1. \score{#2} #3}\par}

\begin{document}

\begin{center}
  {\Large\bfseries 形式语言与编译 2024 考试题（回忆版）}\\
  \vspace{0.3em}
  {\normalsize 来源：\kw{2024形编考试题回忆版.docx}}\\
  {\small 参考：\url{https://xjtu.men/t/topic/12940}}
\end{center}

\textbf{说明：}本文件由回忆版材料整理为 \LaTeX{} 题面。原材料中缺失的数据、选项和图表保持缺失状态，不额外补全为答案。

\qtitle{一}{20}{选择题（10$\times$2）}

回忆版只记录了以下若干题意：
\begin{enumerate}
  \item 下面哪个不能识别正则语言。
  \item 编译的几个步骤，前后端接口。
  \item 一个字母表不能组成哪个字符串。
  \item 递归下降分析表的生成。
  \item 控制流代码翻译，考了 \kw{and}（布尔运算）。
  \item PDA 的状态转移（PDA 有关的选择题不少于 2 道）。
  \item 编译器前后端中间组件。
  \item 什么不属于中间语言（选项里有一个 \kw{sw} 命令，其他是三元、四元）。
\end{enumerate}

\qtitle{二}{15}{将 NFA 转换为最小 DFA 再转 RE}

回忆版未记录具体 NFA 迁移表或状态图。

\qtitle{三}{5}{词法分析}

已给出一个简单的词法分析 DFA。
\begin{enumerate}[label=(\arabic*)]
  \item 根据输入串写出输出。
  \item 判断事实优先级（AAA 和 ADD，即自加运算和加法运算）。
\end{enumerate}

\qtitle{四}{15}{文法修剪}

\begin{enumerate}[label=(\arabic*)]
  \item 修改文法，消除无效符号。
  \item 修改文法，消除左递归。
  \item 求某文法的 FIRST、FOLLOW 集并判断其是否为 LL(1) 文法。
\end{enumerate}

\qtitle{五}{10}{语法分析}

已知句子 \kw{+2++3*(+4)} 和对应的语法树（树已给出）。

文法：
\[
E\gto E*E \mid E+E \mid +i \mid (E).
\]

\begin{enumerate}[label=(\arabic*)]
  \item 对这个句子进行规范规约；写出语法树的直接短语、句柄。
  \item 判断文法是否为歧义，为什么。
\end{enumerate}

\qtitle{六}{10}{写出 itemDFA 并判断冲突}

写出 itemDFA，并判断是否存在冲突；如果存在，写出所有存在的冲突，任意消除其中一个。

文法为：
\[
S'\gto E,\qquad E\gto E+E \mid +E \mid +i.
\]

\qtitle{七}{15}{语义分析}

有程序段：
\begin{lstlisting}
int a[2, 3];

void f(int b(), int k){
  if(k > 0 && k < 3) a[k-1,k] = b(k, 1)
  else print 0;
}
\end{lstlisting}

\begin{enumerate}[label=(\arabic*)]
  \item 对 \kw{f} 函数声明行进行语义分析（自然语言或者符号表均可）。
  \item 对下面语句进行分析，写出三地址/四元式。（未像 2017 题一样给出语义动作）
\end{enumerate}

\begin{lstlisting}
if(k > 0 && k < 3) a[k-1,k] = b(k, 1)
else print 0;
\end{lstlisting}

\qtitle{八}{10}{运行时存储空间组织}

有程序段：
\begin{lstlisting}
int foo(int y;){
  int z;

  void bar(int x; int soo())
  {
    if(x > 3) bar(x/3, soo());
    else z = soo(x);
    print z;
  }

  int row(int x)
  {
    y = x + 5;
    return y;
  }

  bar(y, raw());
}

foo(6);
\end{lstlisting}

注：回忆版中调用处写作 \kw{raw()}，而上方过程名写作 \kw{row}，此处保留原文。

栈结构：
\begin{center}
\begin{tabular}{|c|}
\hline
形参\\
\hline
访问链\\
\hline
控制链\\
\hline
返回地址\\
\hline
变量\\
\hline
\end{tabular}
\end{center}

写出执行到 \kw{return y;} 时候的栈快照。（未像 2017 题一样给出表格，需要自己画。）

\end{document}
